% Nouveau EGA V (cote 25): archive cote 25, site page 170 (PDF page 171).
% typescript. Draft compared: none.
% AI-transcribed by Claude Opus 5.5 (claude-opus-5-5, adaptive thinking, effort high), Message Batches API, batch ega25b (msgbatch_01A7n5YuXubUPNK1zsepi8RG), 2026-11-13; not yet checked by eye.
\scanpage{25/170}
\mspage{\hand{5}}
donner) alors $V^2$ peut se décrire intrinsèquement (i.e. indépendammant de u)
\xout{xxxxx} en termes de l'homomorphisme précédent, comme étant``le sous-préschéma
qui rend l'homomorphisme précédent non dégénéré'' , i.e. qui représente le
foncteur suivant $\emptyset : (\mathrm{Sch})^{o}_{/V^1 - V''} \longrightarrow (\mathrm{Ens})$ \xout{xxx\unsure{F(T)}xxxxxxxxxxxx}, sous-fonc-%
teur du foncteur final, pour lequel $\emptyset(T)\neq \emptyset$ sss l'homomorphisme $\underline{m}_T$ déduit
de $\underline{m}$ par le changement de base $T \to V^1 - V''$ est tel que $\underline{k}_T$ soit inversible et
$\underline{m}_T$ corresponde\xout{xxx} à une fomme quadratique non dégénérée sur $\underline{W}_T$ à valeurs dans
$\underline{k}_T^{-1}$ .

De plus, il y a lieu \xout{xx} d'expliciter l'ensemble $V^o - V'' - V^2$ \xout{xxxxxx}
(\unsure{points} non supersinguliers de $V^o$ rel à Y, et pour la codimension m dans X)
comme étant formé des $x \in V^o$ tels que a) $V^o$ soit plat sur Y en x , et
b) $V^o_y$ admette x comme point géométriquement régulier (i.e. lisse) \uline{ou}
géométriquement \unsure{singulier|siggulier} quadratique ordinaire, \hand{c)}de codimension m dans $X_y$
au point x . (NB y dédigne le point de Y sous x). Dire également qu'en ces
poihts, $V^o \longrightarrow X$ est une immersion transversalement régulière relativement à
Y, de codimension m , (et on aura déjà dit, comme conséquence triviale de
l'égalité des rangs, que quand il y a un $\underline{E}$,u , alors $\underline{E}\otimes\underline{O}_{V^o} \to \underline{N}$ est un
isomorphisme en les poihts en question).

Bien entendu, une fois prouvé les invariances que je viens de signa-%
ler, la définition de $V^1, V'', V^2$ , \xout{xxxdxxxxxxxxaxixx} est claire, indépendamment
de l'existence globale d'un homomorphisme surjectif u:$\underline{E}$ \quad $\underline{J}$ ; d'autre part,
la proposition esquissée fin de la page 16 pour le cas où Y lui-même est
lisse sur un S, où V''= $\emptyset$ et où \inserted{$V^o$} $V^1$ sont lisses sur S avec la \hand{``}bonne\hand{''} dimension
relative sur S, \xout{xixxxxixxxdødixxxxxixx} savoir la caractérisation de $V^2$ comme le
préschéma de ramification de $V^1 \longrightarrow Y$ , prend également un sens intrinsèque,
indépendammant du choix d'un $\underline{E}$.

Pour terminer, on peut expliciter le cas où $V^o$ est défini comme image
inverse, par un \inserted{Y-}morphisme u: $X \to Z$, d'un sous-préschéma $Z^o$ de Z, en suppo-%
\mspage{\hand{170}}
